\documentclass[11pt,a4paper]{article}

\usepackage[utf8]{inputenc}
\usepackage[T1]{fontenc}
\usepackage{geometry}
\geometry{margin=1in}
\usepackage{booktabs}
\usepackage{hyperref}
\usepackage{xcolor}
\usepackage{enumitem}
\usepackage{longtable}

\hypersetup{
    colorlinks=true,
    linkcolor=blue!60!black,
    urlcolor=blue!60!black
}

\title{Market Regime Detection for RL-Based Portfolio Allocation\\[0.4em]
\large Current Project Architecture}
\author{Research Workspace Snapshot}
\date{\today}

\begin{document}
\maketitle
\tableofcontents
\newpage

\section{Purpose}

This repository studies whether weekly latent market regimes improve ETF
allocation across SPY, TLT, GLD, and cash.

The current canonical workflow is:
\begin{enumerate}[nosep]
    \item build weekly price and macro features
    \item fit a Gaussian HMM regime model
    \item infer full-sample causal regime posteriors
    \item add weekly FinBERT news features
    \item train a DQN agent
    \item evaluate against static and heuristic baselines
\end{enumerate}

\section{Canonical Components}

\subsection{Data Acquisition and Feature Engineering}

\begin{itemize}[nosep]
    \item \texttt{scripts/fetch\_yahoo\_seed\_data.py}: fetches daily market data
    \item \texttt{scripts/fetch\_fred\_macro\_panel.py}: fetches FRED macro series
    \item \texttt{scripts/build\_project\_datasets.py}: builds weekly price, macro, and target tables
    \item \texttt{scripts/fetch\_asset\_news.py}: fetches weekly GNews articles
    \item \texttt{scripts/news\_sentiment.py}: scores weekly news with FinBERT
\end{itemize}

The canonical pre-HMM table is:
\begin{verbatim}
data/processed/model_state_weekly_price_macro.csv
\end{verbatim}

The canonical raw news sentiment table is:
\begin{verbatim}
data/raw/news_sentiment/all_assets_news_weekly_finbert.csv
\end{verbatim}

\subsection{HMM Baseline}

The HMM source of truth is:
\begin{verbatim}
scripts/train_hmm_regimes.py
\end{verbatim}

Canonical settings:
\begin{itemize}[nosep]
    \item feature preset: \texttt{regime\_core}
    \item selection mode: \texttt{pipeline}
    \item development window: 2014-07-01 to 2020-12-31
\end{itemize}

The HMM writes development-window artifacts under:
\begin{verbatim}
output/hmm/
\end{verbatim}

\subsection{Merged Notebook Pipeline}

The current notebook workflow is:
\begin{itemize}[nosep]
    \item \texttt{full\_pipeline/01\_hmm\_regime\_pipeline.ipynb}
    \item \texttt{full\_pipeline/02\_rl\_dqn\_with\_hmm\_news.ipynb}
    \item \texttt{full\_pipeline/03\_evaluation\_backtest.ipynb}
\end{itemize}

Shared notebook glue lives in:
\begin{verbatim}
full_pipeline/_pipeline_utils.py
\end{verbatim}

The merged pipeline writes:
\begin{verbatim}
output/full_pipeline/news_features_weekly_finbert_5assets.csv
output/full_pipeline/hmm_regimes_full_sample.csv
output/full_pipeline/model_state_weekly_hmm_news.csv
output/full_pipeline/rl_validation_actions.csv
output/full_pipeline/rl_locked_test_actions.csv
\end{verbatim}

\subsection{Evaluation}

Reusable backtesting and reporting live in:
\begin{verbatim}
evaluation/
\end{verbatim}

This package is the canonical evaluation layer for:
\begin{itemize}[nosep]
    \item default action templates
    \item split-aware datasets
    \item policy backtests
    \item metric tables and equity-curve reporting
\end{itemize}

\section{Execution Flow}

\subsection{Price and Macro Data}

\begin{enumerate}[nosep]
    \item Yahoo daily OHLCV data is fetched.
    \item FRED macro series are fetched.
    \item Weekly market features, weekly macro features, and forward weekly asset targets are built.
    \item These tables are joined into \texttt{model\_state\_weekly\_price\_macro.csv}.
\end{enumerate}

\subsection{News}

\begin{enumerate}[nosep]
    \item Weekly GNews articles are collected for SPY, TLT, GLD, Cash, QQQ, VIX, and TNX.
    \item FinBERT scores article-level sentiment.
    \item The merged notebook pipeline aggregates five weekly sentiment features:
    \texttt{news\_finbert\_compound\_spy},
    \texttt{news\_finbert\_compound\_tlt},
    \texttt{news\_finbert\_compound\_gld},
    \texttt{news\_finbert\_compound\_vix},
    \texttt{news\_finbert\_compound\_tnx}.
\end{enumerate}

\subsection{HMM to RL}

\begin{enumerate}[nosep]
    \item The HMM is fit on the price+macro state using the \texttt{regime\_core} preset.
    \item The selected HMM is used to infer full-sample causal filtered probabilities.
    \item The merged state adds:
    \begin{itemize}[nosep]
        \item \texttt{regime\_filtered}
        \item \texttt{regime\_viterbi}
        \item \texttt{filtered\_prob\_regime\_*}
        \item the five weekly news features
    \end{itemize}
    \item A DQN policy is trained on that merged state.
    \item Validation and locked-test actions are exported for evaluation.
\end{enumerate}

\section{Current Status}

The active HMM baseline is healthy enough for downstream use:
\begin{itemize}[nosep]
    \item 6 of 24 grid-search candidates survive the hard filters
    \item selected project model: \texttt{K=2, n\_pca=8, cov=diag, seed=7}
    \item interpretability: \texttt{8/8}
    \item full-sample regime counts are roughly balanced
\end{itemize}

The merged notebook pipeline has been executed end-to-end in the current
workspace and produces the HMM, RL, and evaluation artifacts listed above.

\section{Historical Components}

The following components remain in the repository but are not the source of truth:

\begin{itemize}[nosep]
    \item \texttt{Attention\_DQN\_Training.ipynb}: exploratory RL notebook
    \item \texttt{pattern\_recognition/}: alternate HMM stack
    \item \texttt{legacy/}: archived Yahoo-era news workflow
    \item \texttt{output/reports/}: historical generated report artifacts
\end{itemize}

\end{document}
